\section{Structure Generation Evaluation}
\label{app:structure-generation-evaluation}

This section details the evaluation protocol for structure generation. We
formalize the crystal solve rate $\mathrm{Sol}_C$~\citep{jin2025oxtal},
including COMPACK~\citep{chisholm2005compack} matching, collision detection,
and experimental-reference selection. We also summarize the differences between
the released OXtal and CLARI evaluators that motivate our re-evaluation of CLARI.

\textbf{Definition of crystal solve rate.}
Let $c\in\{1,\ldots,C\}$ index targets and
$s\in\{1,\ldots,n_c\}$ index generated candidates for target $c$. The
crystal-level solve rate is
\begin{equation}
    \mathrm{Sol}_C
    =\frac{1}{C}\sum_{c=1}^{C}\max_{1\leq s\leq n_c}
        \mathrm{sol}_{c,s},
\end{equation}
where $\mathrm{sol}_{c,s}$ indicates whether candidate $x_{c,s}$ solves target
$c$:
\begin{equation}
    \mathrm{sol}_{c,s}
    =\mathbf{1}\!\left[
        \mathrm{col}_{c,s}=0
        \ \land\ n_{15}(x_{c,s},t^\star_{c,s})\geq 8
        \ \land\ r_{15}(x_{c,s},t^\star_{c,s})<2\,\text{\AA}
    \right],
\end{equation}
where $\mathrm{col}_{c,s}$ is the collision indicator; $n_{15}$ and $r_{15}$
denote the number of matched molecules and RMSD, respectively, returned by
COMPACK at packing size 15; and $t^\star_{c,s}$ is the selected experimental
reference.

Intuitively, a candidate is collision-free ($\mathrm{col}_{c,s}=0$) if all
nonbonded heavy-atom pairs are separated beyond the collision threshold.
Specifically,
\begin{equation}
    \mathrm{col}_{c,s}
    =\mathbf{1}\!\left[
        \exists(i,j):
        r_i^{\mathrm{vdW}}+r_j^{\mathrm{vdW}}-d_{ij}^{(c,s)}
        \geq 0.7\,\text{\AA}
    \right],
\end{equation}
where $d_{ij}^{(c,s)}$ is the distance between a nonbonded heavy-atom pair
$(i,j)$, and $r_i^{\mathrm{vdW}}$ and $r_j^{\mathrm{vdW}}$ are their van der
Waals radii.

For targets with multiple experimental references, we select the best-matching
reference as $t^\star_{c,s}$. Let $T(c)$ denote the reference set for target $c$
and define
\begin{equation}
    P_{c,s}=\{t\in T(c):n_{15}(x_{c,s},t)\geq 8\}.
\end{equation}
OXtal selects
\begin{equation}
    t^\star_{c,s}=
    \begin{cases}
        \displaystyle\arg\min_{t\in P_{c,s}} r_{15}(x_{c,s},t),
        & P_{c,s}\neq\varnothing,\\[3pt]
        \displaystyle\arg\max_{t\in T(c)} n_{15}(x_{c,s},t),
        & P_{c,s}=\varnothing.
    \end{cases}
\end{equation}

\textbf{Differences between the OXtal and CLARI evaluators.}
The released CLARI evaluator differs from the OXtal evaluator in two aspects:
(1) collision detection, including the treatment of intermolecular collisions,
and (2) selection of the best experimental reference for targets with multiple
reference structures. For a consistent comparison, we therefore re-evaluate
CLARI-M, CLARI-L, and CLARI-H using the OXtal definitions.